% GENERATED FILE. Edit canonical lessons, then run npm run generate:books.
% canonical-source-hash: fnv1a64:2d5a26075e43517c
% canonical-lessons: FR-C507-tournoi, FR-C507-medaille, FR-C507-podium, FR-C507-champion, FR-C507-affronter

\chapter{Le tournoi}
\label{ch:fr-le-tournoi}
\hlchaptermodality{507}

\begin{chapteropening}
\textbf{By the end of this chapter:} \emph{I can talk about tournaments and winners.}

Complete the last lesson of chapter 507: I can talk about tournaments and winners.
\end{chapteropening}

\section[le tournoi]{le tournoi — tournament}

\label{lesson:FR-C507-tournoi}



\begin{quote}
Before the new one: say the French for (se qualifier) to qualify, then the French for winner.
\end{quote}

\subsection*{You'll want to know: le tournoi}
\textbf{le tournoi} — \textquotedblleft{}tournament\textquotedblright{}.

\begin{quote}
Notre école organise un tournoi de basket.
\end{quote}

\begin{quote}
\emph{Our school is organising a basketball tournament.}
\end{quote}

In the Middle Ages, a contest between knights.

\begin{cousinweb}
\emph{le tournoi} is built on \emph{tournoyer} \textquotedblleft{}to wheel about\textquotedblright{}, from \emph{tourner} \textquotedblleft{}to turn\textquotedblright{}, which comes from Latin \emph{tornāre} \textquotedblleft{}to turn on a lathe\textquotedblright{} (Lewis \& Short).
\end{cousinweb}

\begin{grammarlens}[title={What you've built}]
One more word for telling, explaining and describing at B1.
\end{grammarlens}

\subsection*{Your turn}
\begin{itemize}
\raggedright
  \item \emph{Say it:} \emph{le tournoi}
  \item \emph{Say it:} \emph{le tournoi}, once more
  \item \emph{Say it:} a sentence of your own with \emph{le tournoi}
  \item \emph{Recall:} say the French for (se qualifier) to qualify, then the French for winner, then say \emph{le tournoi} again
\end{itemize}

\begin{tcolorbox}[breakable,colback=teal!4,colframe=teal!35!black,title={Before you move on}]
What does le tournoi mean? (\textquotedblleft{}Tournament\textquotedblright{}.) Say it once more.
\end{tcolorbox}
\section[la médaille]{la médaille — medal}

\label{lesson:FR-C507-medaille}



\begin{quote}
Before the new one: say the French for winner, then the French for tournament.
\end{quote}

\subsection*{You'll want to know: la médaille}
\textbf{la médaille} — \textquotedblleft{}medal\textquotedblright{}.

\begin{quote}
Elle a gagné une médaille d'argent.
\end{quote}

\begin{quote}
\emph{She won a silver medal.}
\end{quote}

\emph{Le revers de la médaille} — the other side of the coin.

\begin{cousinweb}
\emph{la médaille} was borrowed from Italian \emph{medaglia} (EtymDB), which goes back to Latin \emph{metallum} \textquotedblleft{}metal, a mine\textquotedblright{}, from Greek \emph{métallon}.
\end{cousinweb}

\begin{grammarlens}[title={What you've built}]
One more word for telling, explaining and describing at B1.
\end{grammarlens}

\subsection*{Your turn}
\begin{itemize}
\raggedright
  \item \emph{Say it:} \emph{la médaille}
  \item \emph{Say it:} \emph{la médaille}, once more
  \item \emph{Say it:} a sentence of your own with \emph{la médaille}
  \item \emph{Recall:} say the French for winner, then the French for tournament, then say \emph{la médaille} again
\end{itemize}

\begin{tcolorbox}[breakable,colback=teal!4,colframe=teal!35!black,title={Before you move on}]
What does la médaille mean? (\textquotedblleft{}Medal\textquotedblright{}.) Say it once more.
\end{tcolorbox}
\section[le podium]{le podium — podium}

\label{lesson:FR-C507-podium}



\begin{quote}
Before the new one: say the French for tournament, then the French for medal.
\end{quote}

\subsection*{You'll want to know: le podium}
\textbf{le podium} — \textquotedblleft{}podium\textquotedblright{}.

\begin{quote}
Trois athlètes français sont montés sur le podium.
\end{quote}

\begin{quote}
\emph{Three French athletes went up on the podium.}
\end{quote}

Said \emph{po-dyom}.

\begin{cousinweb}
\emph{le podium} is a learned borrowing from Latin \emph{podium} \textquotedblleft{}a raised platform\textquotedblright{} (EtymDB), from Greek \emph{pódion} \textquotedblleft{}a little foot, a base\textquotedblright{}.
\end{cousinweb}

\begin{grammarlens}[title={What you've built}]
One more word for telling, explaining and describing at B1.
\end{grammarlens}

\subsection*{Your turn}
\begin{itemize}
\raggedright
  \item \emph{Say it:} \emph{le podium}
  \item \emph{Say it:} \emph{le podium}, once more
  \item \emph{Say it:} a sentence of your own with \emph{le podium}
  \item \emph{Recall:} say the French for tournament, then the French for medal, then say \emph{le podium} again
\end{itemize}

\begin{tcolorbox}[breakable,colback=teal!4,colframe=teal!35!black,title={Before you move on}]
What does le podium mean? (\textquotedblleft{}Podium\textquotedblright{}.) Say it once more.
\end{tcolorbox}
\section[le champion]{le champion — champion}

\label{lesson:FR-C507-champion}



\begin{quote}
Before the new one: say the French for medal, then the French for podium.
\end{quote}

\subsection*{You'll want to know: le champion}
\textbf{le champion} — \textquotedblleft{}champion\textquotedblright{}.

\begin{quote}
Il est champion de natation de sa région.
\end{quote}

\begin{quote}
\emph{He is the swimming champion of his region.}
\end{quote}

A woman is \emph{la championne}.

\begin{cousinweb}
\emph{le champion} is inherited from Medieval Latin \emph{campiō} \textquotedblleft{}a fighter in a duel\textquotedblright{} (EtymDB), a Germanic word built on the word for a battlefield, related to Latin \emph{campus} \textquotedblleft{}field\textquotedblright{}.
\end{cousinweb}

\begin{grammarlens}[title={What you've built}]
One more word for telling, explaining and describing at B1.
\end{grammarlens}

\subsection*{Your turn}
\begin{itemize}
\raggedright
  \item \emph{Say it:} \emph{le champion}
  \item \emph{Say it:} \emph{le champion}, once more
  \item \emph{Say it:} a sentence of your own with \emph{le champion}
  \item \emph{Recall:} say the French for medal, then the French for podium, then say \emph{le champion} again
\end{itemize}

\begin{tcolorbox}[breakable,colback=teal!4,colframe=teal!35!black,title={Before you move on}]
What does le champion mean? (\textquotedblleft{}Champion\textquotedblright{}.) Say it once more.
\end{tcolorbox}
\section[affronter]{affronter — to face, to take on}

\label{lesson:FR-C507-affronter}



\begin{quote}
Before the new one: say the French for podium, then the French for champion.
\end{quote}

\subsection*{You'll want to know: affronter}
\textbf{affronter} — \textquotedblleft{}to face, to take on\textquotedblright{}.

\begin{quote}
Demain, notre équipe affronte les champions.
\end{quote}

\begin{quote}
\emph{Tomorrow our team takes on the champions.}
\end{quote}

\emph{S'affronter} — to face each other.

\begin{cousinweb}
\emph{affronter} comes from Old French \emph{afronter} \textquotedblleft{}to face\textquotedblright{} (EtymDB), built on \emph{front} \textquotedblleft{}forehead\textquotedblright{}, from Latin \emph{frōns} (Lewis \& Short): to stand forehead to forehead.
\end{cousinweb}

\begin{grammarlens}[title={What you've built}]
One more word for telling, explaining and describing at B1.
\end{grammarlens}

\subsection*{Your turn}
\begin{itemize}
\raggedright
  \item \emph{Say it:} \emph{affronter}
  \item \emph{Say it:} \emph{affronter}, once more
  \item \emph{Say it:} a sentence of your own with \emph{affronter}
  \item \emph{Recall:} say the French for podium, then the French for champion, then say \emph{affronter} again
\end{itemize}

\begin{tcolorbox}[breakable,colback=teal!4,colframe=teal!35!black,title={Before you move on}]
What does affronter mean? (\textquotedblleft{}To face, to take on\textquotedblright{}.) Say it once more.
\end{tcolorbox}
